\documentclass[11pt]{article}
\usepackage[utf8]{inputenc}
\usepackage[T1]{fontenc}
\usepackage[margin=1in]{geometry}
\usepackage{hyperref}
\usepackage{enumitem}
\hypersetup{colorlinks=true,linkcolor=black,urlcolor=blue,citecolor=blue}
\title{Robot Logic and Conditional Coding}
\author{Piter Garcia}
\date{March 20, 2026}
\begin{document}
\maketitle
\section*{Packet Use}
This printable lesson packet is generated from the paired Markdown guide. The Markdown version remains the clickable, neurodivergent-friendly working copy; this PDF carries the same lesson substance in a formal handout format. Evidence claims in this packet are based on transcript-derived lesson notes and the cleaned transcript files in this workspace.

\section*{Quick Scan}
\begin{itemize}[leftmargin=*]
\item \textbf{Grade:} Grade 5
\item \textbf{Grade Hub:} Notes [../grade-5-Notes.md]
\item \textbf{Schedule Reference:} Spring2026\_TeachingSchedule\_Derek.md [../../school/Spring2026\_TeachingSchedule\_Derek.md]
\item \textbf{Workbook Sheet:} \texttt{5th} in \texttt{Teaching\_Placement-IGNITE Curriculum\_ GCSD25-26.xlsx}
\item \textbf{Lesson Focus:} Students move from basic robot control to conditional logic, using sensors, events, or logic choices to change what the robot does.
\item \textbf{CS Alignment Strength:} Strong CS lesson
\item \textbf{LaTeX Source:} robot-logic-and-conditional-coding.tex [./robot-logic-and-conditional-coding.tex]
\item \textbf{Bib File:} robot-logic-and-conditional-coding.bib [./robot-logic-and-conditional-coding.bib]
\item \textbf{Artifacts Folder:} artifacts/ [./artifacts/]
\end{itemize}

\section*{Lesson Identity}
\begin{itemize}[leftmargin=*]
\item \textbf{Likely Class Blocks:} Day 2 -- 10:45-11:35 -- Pickett, Day 3 -- 10:45-11:35 -- McGlashon, Day 5 -- 10:45-11:35 -- Pecora
\item \textbf{Main Lesson Family:} Robot Logic and Conditional Coding
\item \textbf{Primary Evidence Base:} Transcript-derived notes plus source transcripts from this teaching-placement workspace.
\item \textbf{Primary External Source Type:} Official curriculum/provider materials.
\end{itemize}

\section*{What We Are Teaching}
\begin{itemize}[leftmargin=*]
\item what a condition means in code
\item how robot behavior changes when a condition is true or false
\item how to compare a simple sequence with a logic-driven program
\item how to debug conditional branches and trigger conditions
\end{itemize}

\section*{CS Standards Alignment}
\begin{itemize}[leftmargin=*]
\item \textbf{1B-AP-10}: Create programs that include sequences, events, loops, and conditionals. Why it fits here: The lesson explicitly targets conditional logic in robot programming.
\item \textbf{1B-AP-11}: Decompose problems into smaller, manageable subproblems. Why it fits here: Students separate condition, action, and test case while planning the task.
\item \textbf{1B-AP-15}: Test and debug a program or algorithm to ensure it runs as intended. Why it fits here: Students must test both branches or outcomes of the condition.
\end{itemize}

\section*{NYS / Local Curriculum Alignment}
\begin{itemize}[leftmargin=*]
\item Aligned to the local Grade 5 robotics sequence in the workbook and schedule.
\item Supports New York-facing problem solving and computational reasoning through branching logic and debugging.
\item This lesson is a strong CS packet because students are no longer only sequencing movement; they are reasoning about logic and alternate outcomes.
\end{itemize}

\section*{Lesson Content and Concepts}
\begin{itemize}[leftmargin=*]
\item conditional
\item branching
\item sensor input
\item test case
\item debugging
\end{itemize}

\section*{Evidence / Proof of Instruction}
\begin{itemize}[leftmargin=*]
\item \textbf{Transcript-derived note:} 260227 Grade 4 5 Shared Conditional Sphero [../../transcript-derived-lessons/260227-grade-4-5-shared-conditional-sphero.md] The shared Grade 4-5 note explicitly frames the work as conditional coding with Sphero robots.
\item \textbf{Transcript-derived note:} 260313 Mixed Shared Sphero Minecraft [../../transcript-derived-lessons/260313-mixed-shared-sphero-minecraft.md] Later mixed notes show students working across robot logic and Minecraft-style command thinking, reinforcing the logic focus.
\item \textbf{Source transcript:} 260227 Grade 4 5 Shared Conditional Sphero [../../transcripts/260227-grade-4-5-shared-conditional-sphero.md] The source transcript supports that conditions and robot responses were taught explicitly, not inferred later.
\item \textbf{Likely student proof:} Robot runs that behave differently based on a sensor or condition, plus explained fixes, count as direct evidence.
\end{itemize}

\section*{Assessment Evidence}
\begin{itemize}[leftmargin=*]
\item Student can explain what condition the robot is checking.
\item Student can build or modify a conditional block sequence.
\item Student tests and explains at least one failure or false trigger.
\end{itemize}

\section*{UDL and Inclusion Supports}
\subsection*{Multiple Representation}
\begin{itemize}[leftmargin=*]
\item Show condition language with if/when visuals and concrete trigger examples.
\item Use a two-column chart: if this happens / then robot does this.
\item Demonstrate one true-case and one false-case outcome.
\end{itemize}

\subsection*{Multiple Action / Expression}
\begin{itemize}[leftmargin=*]
\item Allow students to predict with cards before coding.
\item Provide prewritten conditional sentence frames.
\item Use paired testing roles so one student watches the condition while the other launches code.
\end{itemize}

\subsection*{Multiple Engagement}
\begin{itemize}[leftmargin=*]
\item Use visible cause-and-effect robot behavior to make logic concrete.
\item Keep tasks framed as puzzles with more than one test round.
\item Invite students to guess what the robot will do before it runs.
\end{itemize}

\section*{How We Are Already Making It Inclusive}
\begin{itemize}[leftmargin=*]
\item The transcript trail shows that students were already beyond basic movement coding and into more advanced reasoning.
\item Robot tasks provide embodied feedback, which supports students who need visible evidence of logic.
\item The lesson already appears to use partner structures and repeated testing cycles.
\end{itemize}

\section*{How to Improve Inclusion Next Time}
\begin{itemize}[leftmargin=*]
\item Add a branch organizer so students record both possible outcomes before launch.
\item Create one simplified conditional task and one enrichment task in the same packet.
\item Capture a short video of a working true/false branch in \texttt{artifacts/}.
\item Add sentence stems for students explaining why the robot chose one action instead of another.
\end{itemize}

\section*{Materials and Setup}
\begin{itemize}[leftmargin=*]
\item Robots and charging setup
\item Student devices
\item Conditional challenge cards
\item Sensor trigger objects or controlled environment
\end{itemize}

\section*{Teacher Moves / Script / Routines}
\begin{itemize}[leftmargin=*]
\item Model a non-conditional program first, then show what changes with a condition.
\item Ask students to test both outcomes instead of only the successful one.
\item Use debug talk that distinguishes wrong code from wrong trigger conditions.
\item Close with a comparison between sequence-only and condition-based solutions.
\end{itemize}

\section*{Common Errors and Debugging}
\begin{itemize}[leftmargin=*]
\item Students may think a condition is just another step, not a decision point.
\item Sensors can create inconsistent results if the environment is uncontrolled.
\item Students may test only one branch and assume the program works.
\end{itemize}

\section*{Extensions and Simplifications}
\begin{itemize}[leftmargin=*]
\item Extension: chain more than one condition or compare two conditional solutions.
\item Simplification: use one obvious trigger and one obvious response.
\item Extension: ask students to explain the logic in plain language before showing the blocks.
\end{itemize}

\section*{Related Local Materials}
\begin{itemize}[leftmargin=*]
\item No direct local lesson packet linked yet.
\end{itemize}

\section*{Artifacts Checklist}
\begin{itemize}[leftmargin=*]
\item Compiled lesson PDF in \texttt{artifacts/}
\item At least one screenshot, student example, or setup photo
\item One short assessment or observation record
\item Updated standards/source bibliography once the \texttt{.bib} file is revised
\item Any teacher reflection or revision notes after reteaching
\end{itemize}

\section*{Selected Official Sources}
Selected official sources that informed this packet are listed below.
ocite{*}

\bibliographystyle{plain}
\bibliography{robot-logic-and-conditional-coding}
\end{document}
